\section{Guards, extensionality, and precision contracts}
\label{sec:guards}

\subsection{Totalizing partial operations}

Primitive-recursive function symbols are total, while mathematical operations
such as division, normalization, chart inversion, and pivoting are partial.
The sound pattern is to define an arbitrary total extension and state all
correctness theorems under a numerical guard.

For reciprocal, a guard may be
\[
  |x|\ge 2^{-g}.
\]
The reciprocal evaluator reads the input at a precision determined by \(g\)
and the requested output precision.  Its object theorem has the form
\[
 \mathsf{Apart}(x,g)\wedge\mathsf{Name}(a,x)
 \to \mathsf{Name}(\mathsf{inv}(a,g),x^{-1}).
\]
The total function's behavior outside the guard is irrelevant.

A mere proposition \(x\ne0\) does not generally yield such a rational margin in
constructive arithmetic.  The margin must be supplied, searched for in a
bounded exact presentation, or retained as an explicit producer obligation.

\subsection{Overlapping branches}

Hard branch decisions are often unnecessary.  Suppose branch \(L\) is sound
when \(x<\beta\), branch \(R\) is sound when \(x>\alpha\), and
\(\alpha<\beta\).  In the overlap \(\alpha<x<\beta\), either branch proves the
same output contract.  A finite-precision comparison can choose one branch
without deciding an exact boundary case.

\begin{principle}[Robust branching]
Use overlapping guards whenever the conclusion is stable on the overlap.  Use
an exact branch label only when a positive rational gap proves that one branch
is uniquely valid.
\end{principle}

This principle covers interval subdivision, coordinate charts, matrix pivots,
sign-dependent formulas, and approximate root labels.

\subsection{Observational equivalence}

Let \(p\approx_n q\) mean that two codes agree to observation precision \(n\).
For regular rational names this can be a dyadic distance bound; for compact
sets it can be finite Hausdorff distance; for root data it can be a matching or
cluster relation; for finite probability laws it can be a rational test metric.

A certified operation \(F\) carries a demand transformer \(D_F\) and proves
\[
 p\approx_{D_F(n)}q
 \quad\Longrightarrow\quad
 F(p)\approx_n F(q).
\]
This is the extensionality theorem.  It replaces quotient construction by a
family of finite relations that is directly usable in iPRA.

\begin{theorem}[Precision-contract composition]
\label{thm:precision-composition}
If \(F\) is extensional with demand \(D_F\) and \(G\) is extensional with
demand \(D_G\), then \(G\circ F\) is extensional with demand
\(D_F\circ D_G\).  The same holds for relational outputs when the component
relations satisfy finite back-and-forth totality.
\end{theorem}

\begin{proof}
From \(p\approx_{D_F(D_G(n))}q\), extensionality of \(F\) gives
\(F(p)\approx_{D_G(n)}F(q)\).  Extensionality of \(G\) gives the desired
output relation.  In the object calculus the proof is two implication
eliminations after substituting the requested precision terms.  For relational
outputs, choose bounded witnesses supplied by the two back-and-forth clauses
and compose their certificates by finite search.
\end{proof}

\subsection{Error ledgers}

A precision contract should report not only a demand but a numerical error
budget.  For example, if
\[
 |x-x'|\le\epsilon_x,
 \qquad |y-y'|\le\epsilon_y,
 \qquad |x|,|x'|\le M_x,
 \qquad |y|,|y'|\le M_y,
\]
then
\[
 |xy-x'y'|\le M_x\epsilon_y+M_y\epsilon_x.
\]
All quantities can be coded rationals, and cross multiplication removes
ordered-field division under positive denominator guards.  Finite sums and
products use accumulators whose local step is a polynomial inequality.

The recommended module exposes each arithmetic rule in three forms:
\begin{enumerate}[label=(\roman*)]
\item an ambient exact rational lemma for testing;
\item a fixed object-language formula over \(\PTm\); and
\item a derivation constructor in \(\Prf\,\PRAAx\).
\end{enumerate}
Only the third form can close an application packet.

\subsection{Regular names without quotient reals}

A regular name can be represented by a fixed evaluator \(a(p,n)\) and the
contract
\[
 |a(p,n)-a(p,m)|\le2^{-n}+2^{-m}.
\]
Arithmetic is formulated as transformations of such contracts.  There is no
need to define a completed real-number quotient inside the object theory.
Equality of represented reals is the observational relation
\[
 \forall n\quad |a(p,n)-a(q,n)|\le 2^{-n+c}
\]
for a fixed harmless shift \(c\).  Every theorem ultimately quantifies over
natural codes and rational inequalities.
